\section{Introduction}
\label{sec:intro}

Post-hoc explanation asks what aspects of an input matter to a model’s prediction. A natural way to answer this question is contrastively: compare the prediction with what would happen if a feature, region, concept, or context were changed. Perturbation-based and counterfactual methods instantiate this comparison by constructing alternative inputs and measuring the resulting change in model output~\citep{fong2017meaningful,dhurandhar2018explanations,goyal2019counterfactual,covert2021explaining}. This has the basic form of an effect measurement: change a factor, compare two responses, and summarize the difference. In practice, that response is often reduced to a scalar magnitude.

A magnitude answers an important question: \emph{how much did the model react?} It leaves a second question open: \emph{what did the reaction mean?} Consider an input and a paired counterfactual that changes one factor. Their prediction difference measures how strongly the model responds to that change. This endpoint comparison, however, shows only the net effect after inputs are fully observed. It does not show how the contrast emerges as information becomes available. We therefore consider a paired reveal: the paired inputs are progressively revealed along matched trajectories, so that their prediction difference can be observed at the same level of revealed information.

As the pair is revealed from an uninformative state, the contrast can evolve in different ways. It may align with the final effect, oppose it, or become large even when the fully observed pair differs little. These distinct behaviors can have identical absolute responses. Treating these responses as equivalent therefore hides behavior that matters for interpretation.

\begin{figure*}[t]
\centering
\begin{minipage}[t]{0.6\linewidth}
\begin{figure}[H]
  \vspace{0pt}
  \centering
  \includegraphics[width=\linewidth]{figures/conceptual/fig1_decaf_overview.pdf}
  \caption{\textbf{\decaf at a glance.} The paired reveal measures a signed stage response, while the clean endpoint supplies an activity gate and orientation. \decaf routes the same response into endpoint-aligned evidence, endpoint-opposed contradiction, or endpoint-null fragility, with exact conservation. The compact summary reports the two headline consequences: mechanism recovery after ordinary magnitude is matched, and substantially more stable cross-protocol diagnostics.}
\label{fig:decaf-overview}
\end{figure}
\end{minipage}%
\hfill
\begin{minipage}[t]{0.38\linewidth}
We introduce \decaf in \figurename~\ref{fig:decaf-overview}, short for \emph{\textsc{D}ecomposition of \textsc{E}vidence, \textsc{C}ontradiction, \textsc{A}nd \textsc{F}ragility}. \decaf resolves this ambiguity by using the final response between the fully observed pair as a semantic reference for the entire reveal process. Its magnitude determines whether the factor has a meaningful final effect, while its sign provides the reference direction. Intermediate responses that agree with this direction are routed to \emph{evidence}, while responses that oppose it are routed to \emph{contradiction}. If the fully observed pair shows little response, there is no meaningful final effect to align with, and intermediate responses are routed to \emph{fragility}. We call this \emph{semantic routing}: the observed responses are assigned different roles according to their relation to the final contrast, without changing the trajectories themselves. The routing is lossless, so evidence, contradiction, and fragility sum exactly to the ordinary response magnitude.
  \textbf{\textbf{\decaf} is model-agnostic}, therefore explaining any black-box model that returns a scalar score.
\end{minipage}

\end{figure*}


\textbf{Relation to prior work.}
Gradient and path methods assign contributions to input coordinates or internal units~\citep{simonyan2014deep,selvaraju2017gradcam,sundararajan2017axiomatic,shrikumar2017deeplift,bach2015lrp}. Removal and sampling methods measure prediction changes under masking or replacement~\citep{fong2017meaningful,fong2019extremal,petsiuk2018rise,covert2021removing}. Counterfactual and causal approaches define meaningful contrasts~\citep{goyal2019counterfactual,chattopadhyay2019causal,janzing2020causal}. Evaluation work shows that baselines, removal operators, and off-manifold inputs can alter explanation behavior~\citep{adebayo2018sanity,hooker2019benchmark,rong2022consistent,jethani2023labelleakage}. \decaf begins after a paired response has been measured. It refines the response by its relation to the clean endpoint. Unlike positive--negative feature attribution, it also contains an endpoint-null branch. Unlike amortized explainers such as FastSHAP, it requires no learned explanation model~\citep{jethani2022fastshap}. A detailed comparison with feature attribution, counterfactual explanation, evaluation frameworks, and efficient black-box methods appears in Appendix~\ref{app:related}.

\textbf{We make five contributions}:
\begin{itemize}[leftmargin=1.4em,itemsep=1pt,topsep=0pt]
    \item \emph{Problem.}
    We identify a semantic gap in perturbation explanation: response magnitude tells us how strongly a model reacts, but not what that reaction means.

    \item \emph{Observation interface.}
    We introduce paired reveal, which tracks a factual--counterfactual contrast along matched trajectories as information becomes available.

    \item \emph{Algorithm.}
    We introduce \decaf, which uses the final contrast to route intermediate responses into evidence, contradiction, and fragility.
    The decomposition is lossless and unique under simple axioms, and adds no model queries beyond the paired trajectories.

    \item \emph{Behavioral meaning.}
    Controlled vision and tabular experiments show that the three components track learned reliance, effect reversal, and fragility across neural, linear, and tree-based models.

    \item \emph{Practical consequence.}
    On ImageNet-9, \decaf recovers response semantics after magnitude is matched and diagnoses what changes across perturbation paths.
    It also outperforms the tested general-purpose attribution baselines on FunnyBirds and ImageNet-1k and scales efficiently to a large DINOv2 vision transformer.

\end{itemize}